% ============================================================
\part{پایه‌های محاسباتی}
% ============================================================

\chapter{الگوریتم چیست؟}
\label{chap:algorithm}

\section{جایگاه در نقشه}

لایه صفر: قبل از کوانتوم. بدون این تعریف، «الگوریتم کوانتومی» با سخت‌افزار یا با یک گیت خاص قاطی می‌شود. در نقشه فصل~\ref{chap:map} هنوز به هامیلتونی نرسیده‌ایم.

\section{ایده شهودی}

الگوریتم یک دستورالعمل محاسباتی است: ورودی مشخص می‌گیرد، با تعداد متناهی گامِ بدون ابهام پیش می‌رود، و خروجی می‌دهد. رایانه، آزمایشگاه، یا کاغذ و قلم مجری‌اند، نه خودِ الگوریتم.

سه پرسش همیشه همراه الگوریتم‌اند: \textbf{درستی} (آیا خروجی همان چیزی است که مسئله می‌خواهد؟)، \textbf{پایان‌پذیری} (آیا حتماً تمام می‌شود؟)، و \textbf{هزینه} (با اندازه ورودی، زمان و حافظه چطور رشد می‌کنند؟).

شهود مفید: دستور پخت غذا الگوریتم نیست اگر «کمی نمک» مبهم بماند. الگوریتم وقتی الگوریتم است که اجراکننده‌های مختلف، با همان ورودی، به همان نتیجه برسند (تا حد خطای اعلام‌شده).

\section{حداقل فرمالیسم}

یک الگوریتم $A$ را می‌توان نگاشتی از نمونه مسئله $x$ به خروجی $A(x)$ دانست، به‌همراه مدلی از هزینه (تعداد گام، عمق مدار، تعداد پرسش از یک اوراکل). در این سند دو تمایز کافی است:
\begin{itemize}
	\item \textbf{مسئله}: چه چیزی پرسیده می‌شود؟ (مثلاً انرژی پایه یک مدل)
	\item \textbf{الگوریتم}: با چه روالی به پاسخ می‌رسیم؟
\end{itemize}
پیچیدگی را اینجا وارد جزئیات نمی‌کنیم. «کوانتومی بودن» مدل اجرا را عوض می‌کند، نه تعریف الگوریتم را.

\section{چه وقت استفاده می‌شود}

وقتی وسوسه می‌شوید \lr{Trotter} یا \lr{VQE} را «خودِ مسئله» بنامید. آن‌ها الگوریتم یا زیرروال‌اند؛ مسئله چیز دیگری است (مثلاً تحول یک مدل هابارد).

\begin{remark}[جای توسعه]
مدل‌های هزینه (مدار، پرسمان، فضا)؛ تمایز الگوریتم دقیق، تقریبی، تصادفی و اکتشافی؛ اشاره به Church--Turing و نسخه کوانتومی آن.
\end{remark}

\chapter{خانواده‌های کلاسیکِ مرتبط}
\label{chap:classical}

\section{جایگاه در نقشه}

هنوز در لایه کلاسیک هستیم. هدف، کاتالوگ الگوریتم‌های عمومی نیست؛ فقط خانواده‌هایی که بعداً با شبیه‌سازی یا الگوریتم کوانتومی تلاقی مفهومی دارند.

\section{ایده شهودی}

هر بار که یک الگوریتم کوانتومی می‌بینید، معمولاً جانشین یا رقیب یک روال کلاسیک است. اگر آن رقیب را نشناسید، نمی‌فهمید کوانتوم دقیقاً چه چیزی را عوض کرده: مدل هزینه را، نمایش حالت را، یا خودِ سؤال را. پنج خانواده برای این سند کافی‌اند.

\paragraph{جستجو و بهینه‌سازی.}
جستجوی ناساخت‌یافته و بهینه‌سازی ترکیبیاتی. در کوانتوم: \lr{Grover}، \lr{QAOA}، و بازپخت کوانتومی.\footnote{\lr{quantum annealing}.}

\paragraph{تبدیل فوریه.}
از حوزه زمان/مکان به بسامد. نسخه کوانتومی آن \lr{QFT} است و ستون برآورد فاز.

\paragraph{دستگاه خطی و زیرفضای \lr{Krylov}.}
حل $Ax=b$ و استخراج اطلاعات طیفی با تکرارهای ماتریسی. در کوانتوم: \lr{HHL} و به‌طور کلی‌تر توابعی از عملگر از راه بلوک‌کدگذاری.\footnote{\lr{block-encoding}.} شبیه‌سازی هامیلتونی هم از همین جنس است: اعمال تابعی از $H$.

\paragraph{نمونه‌برداری و مونت‌کارلو.}
برآورد یک مقدار با نمونه‌های تصادفی. نسخه کوانتومی اغلب برآورد دامنه است. در فیزیک بس‌ذره‌ای،\footnote{\lr{many-body}.} مونت‌کارلو کوانتومی خودش یک رقیب کلاسیک جدی است.

\paragraph{شبیه‌سازی کلاسیک سامانه‌های کوانتومی.}
این مهم‌ترین رقیب مستقیم پژوهش شماست: قطری‌سازی دقیق و شبیه‌سازی بردار حالت\footnote{\lr{state vector simulation}.} (دقیق، نمایی در تعداد ذرات)؛ مونت‌کارلو کوانتومی (گاهی مقیاس بهتر، با محدودیت علامت)؛ و شبکه تانسوری\footnote{\lr{tensor network}.} و حالت حاصل‌ضرب ماتریسی\footnote{\lr{matrix product state}.} از جمله \lr{DMRG}، که برای درهم‌تنیدگی محدود قوی‌اند~\cite{Verstraete2008,Cirac2021}. الگوریتم کوانتومی وقتی معنا دارد که یکی از این مسیرها از نظر مقیاس، دقت، یا نوع سؤال ناکافی شود.

\section{حداقل فرمالیسم}

هزینه کلاسیک شبیه‌سازی یک سامانه کوانتومی عمومی معمولاً با بُعد فضای هیلبرت ( یعنی به‌طور نمایی با تعداد ذرات )رشد می‌کند. استثناها همان روش‌هایی‌اند که ساختار (محلی بودن، نبود علامت، درهم‌تنیدگی کم) را استثمار می‌کنند.


\section{چه وقت استفاده می‌شود}

قبل از انتخاب یک الگوریتم کوانتومی بپرسید: همین سؤال را کلاسیک تا چه اندازه و با چه روشی می‌توانم جواب بدهم؟

\begin{remark}[جای توسعه]
جدول مقیاس‌بندی کیفی میان قطری‌سازی، \lr{MPS}، \lr{QMC} و مدار کوانتومی؛ قانون مقیاس‌بندی هزینه در هر خانواده.
\end{remark}

\chapter{مدل محاسباتی کوانتومی}
\label{chap:qmodel}

\section{جایگاه در نقشه}

این فصل می‌گوید رایانه کوانتومی و الگوریتم کوانتومی \emph{کجا} روی نقشه می‌نشینند: یک مدل اجرای تازه برای همان مفهوم الگوریتم، نه صرفاً رایانه‌ای سریع‌تر.

\section{ایده شهودی}

رایانه کلاسیک بیت را در $0$ یا $1$ نگه می‌دارد. رایانه کوانتومی با حالت‌هایی در فضای هیلبرت کار می‌کند. سه منبع محاسبات را از کلاسیک جدا می‌کنند: برهم‌نهی،\footnote{\lr{superposition}.} تداخل،\footnote{\lr{interference}.} و درهم‌تنیدگی.\footnote{\lr{entanglement}.}

الگوریتم کوانتومی روالی است که از این منابع برای یک مسئله با ورودی کلاسیک یا کوانتومی استفاده می‌کند و در پایان با اندازه‌گیری به خروجی کلاسیک (یا به یک حالت کوانتومی مفید) می‌رسد~\cite{NielsenChuang2010}. شبیه‌سازی کوانتومی یک کلاس از مسئله‌ها روی همین مدل است، نه تعریف خودِ مدل~\cite{Feynman1982,Lloyd1996,Georgescu2014}.

سخت‌افزار سه شکل اصلی دارد که بعداً تکرار می‌شوند:
\begin{enumerate}
	\item \textbf{مدار گیتی (دیجیتال):} دنباله گیت‌ها از یک مجموعه گیت جهانی؛\footnote{\lr{universal gate set}.}
	\item \textbf{آنالوگ:} مهندسی خودِ هامیلتونی روی دستگاه؛
	\item \textbf{آدیاباتیک:} محاسبه با مسیر کند $H(s)$.
\end{enumerate}
این سه، سه رهیافت اجرا هستند. یک هدف پژوهشی (مثلاً انرژی پایه هابارد) ممکن است با هر سه دنبال شود.

\section{حداقل فرمالیسم}

مدل مدار: حالت $\lvert\psi\rangle$ روی $n$ کیوبیت، تحول یکانی $U$ ساخته‌شده از گیت‌های محلی، سپس اندازه‌گیری در یک پایه. خروجی الگوریتم نوعاً یک برآورد از $\langle\psi\rvert O\lvert\psi\rangle$ است، نه کل تابع موج مگر در مسائل کوچک.

تمایز کوتاه: گیت\footnote{\lr{gate}.} دستور موضعی است؛ مدار الگوریتم در مدل دیجیتال است؛ زمان‌بندی هامیلتونی الگوریتم در مدل آنالوگ یا آدیاباتیک است.

\section{ارتباط}

الگوریتم به‌معنای عام در فصل~\ref{chap:algorithm}، سه لایه در فصل~\ref{chap:three-layers}، دیجیتال در برابر آنالوگ در فصل~\ref{chap:hardware}، و تحول در برابر آدیاباتیک در فصل~\ref{chap:evo-adi} آمده است.

\section{چه وقت استفاده می‌شود}

وقتی می‌پرسید الگوریتم‌های کوانتومی چه هستند، و شبیه‌سازی و رایانه کوانتومی کجای این بحث‌اند. پاسخ کوتاه: رایانه مدل اجراست؛ الگوریتم روال است؛ شبیه‌سازی یک خانواده از مسائل روی همان مدل است.

\begin{remark}[جای توسعه]
مدل پرسمان و \lr{BQP} در حد یک پاراگراف؛ نویز و ناهمدوسی به‌عنوان فاصله مدل آرمانی از دستگاه؛ \lr{early-FTQC} در برابر \lr{NISQ}.
\end{remark}

\chapter{سه لایه: الگوریتم، شبیه‌سازی، ابزار}
\label{chap:three-layers}

\section{جایگاه در نقشه}

این فصل ستون مفهومی کل سند است. درهم‌ریختگی پژوهشی معمولاً از قاطی کردن این سه لایه می‌آید. جدول عملی «اگر می‌خواهید\ldots» در فصل~\ref{chap:map} است.

\section{ایده شهودی}

سه مفهوم نزدیک‌اند ولی یکی نیستند. شکل~\ref{fig:three-layers} ترتیب منطقی‌شان را نشان می‌دهد.

\begin{figure}[htbp]
	\centering
	\begin{latin}
			\setstretch{1}
			% سبک مشترک نمودارهای سند پایه‌ها (فقط داخل محیط latin فراخوانی شود)
\tikzset{
	qafbox/.style={
		draw=black!55,
		rounded corners=2pt,
		align=center,
		inner sep=3pt,
		thick,
		fill=white,
		font=\footnotesize\linespread{1}\selectfont,
		minimum height=0.95cm
	},
	qafwide/.style={
		qafbox,
		text width=2.55cm
	},
	qafnarr/.style={
		qafbox,
		text width=2.05cm
	},
	qaftitle/.style={
		font=\small\bfseries\linespread{1}\selectfont,
		align=center
	},
	qafarr/.style={
		-{Latex[length=2.2mm]},
		thick,
		draw=black!65
	},
	qafpanel/.style={
		draw,
		thick,
		rounded corners=4pt,
		inner xsep=8pt,
		inner ysep=8pt
	}
}

\begin{tikzpicture}[
	font=\footnotesize\linespread{1}\selectfont,
	node distance=0.55cm
	]
	\node[qafbox, text width=7.6cm, fill=blue!10] (l2) {%
		\textbf{Layer 2 --- Quantum simulation}\\[2pt]
		Which physical model is being represented?};
	\node[qafbox, text width=7.6cm, fill=green!10, below=of l2] (l1) {%
		\textbf{Layer 1 --- Quantum algorithm}\\[2pt]
		Which procedure solves that question?};
	\node[qafbox, text width=7.6cm, fill=orange!10, below=of l1] (l3) {%
		\textbf{Layer 3 --- Implementation tools}\\[2pt]
		Trotter, LCU, QSP, QPE, VQE, adiabatic path};
	\draw[qafarr] (l2) -- (l1);
	\draw[qafarr] (l1) -- (l3);
\end{tikzpicture}

		\end{latin}
	\caption{سه لایه نزدیک ولی جدا: شبیه‌سازی می‌گوید چه مدلی مطالعه می‌شود؛ الگوریتم می‌گوید با چه روالی؛ ابزار می‌گوید آن روال چطور ساخته می‌شود.}
	\label{fig:three-layers}
\end{figure}

\paragraph{لایه ۱ --- الگوریتم کوانتومی.}
دستورالعمل محاسباتی برای حل یک مسئله. مثال: با این مدار، انرژی پایه را تا خطای $\varepsilon$ برآورد کن.

\paragraph{لایه ۲ --- شبیه‌سازی کوانتومی.}
استفاده از یک سامانه کوانتومی برای بازنمایی و مطالعه یک سامانه یا مدل کوانتومی دیگر. اینجا هدف پژوهشی است: مدل هابارد، یک مولکول، یک نظریه پیمانه‌ای شبکه‌ای.\footnote{\lr{lattice gauge theory}.} شبیه‌سازی می‌تواند دیجیتال، آنالوگ، یا ترکیبی باشد. شبیه‌سازی یک گیت خاص نیست~\cite{CiracZoller2012,Georgescu2014}.

\paragraph{لایه ۳ --- روش‌های پیاده‌سازی.}
ابزارهایی برای تحقق یک هدف محاسباتی مشخص: \lr{Trotter}، ترکیب خطی یکانی‌ها (\lr{LCU})، پردازش سیگنال کوانتومی (\lr{QSP})، برآورد فاز (\lr{QPE})، \lr{VQE}، و تحول آدیاباتیک.

جمله جهت‌یاب: شبیه‌سازی می‌گوید چه مدلی را مطالعه می‌کنیم؛ الگوریتم می‌گوید با چه روالی؛ ابزار می‌گوید آن روال را چطور روی دستگاه می‌سازیم.

\section{حداقل فرمالیسم}

\begin{table}[htbp]
	\centering
	\footnotesize
	\setstretch{1.15}
	\caption{نمونه جداسازی سه لایه.}
	\label{tab:three-layers}
	\begin{tabularx}{\textwidth}{@{}>{\raggedright\arraybackslash}p{4.4cm} c >{\raggedright\arraybackslash}X@{}}
			\toprule
			\rowcolor{gray!14}
			\textbf{می‌خواهید} & \textbf{لایه} & \textbf{نمونه} \\
			\midrule
			مسئله را نام ببرید
			& ۲
			& دینامیک هابارد پس از تغییر ناگهانی هامیلتونی (\lr{quench})
			\\
			\rowcolor{gray!8}
			روال حل را بنویسید
			& ۱
			& آماده‌سازی $\lvert\psi_0\rangle$، اعمال $U(t)$، اندازه‌گیری $O$
			\\
			$U(t)$ را بسازید
			& ۳
			& فرمول حاصل‌ضربی یا \lr{QSP}
			\\
			\rowcolor{gray!8}
			ویژه مقدار را بخوانید
			& ۱ و ۳
			& \lr{QPE} که زیرروالش شبیه‌سازی کنترل‌شده است
			\\
			حالت پایه روی \lr{NISQ}
			& ۱ و ۳
			& \lr{VQE}: آنزات وردشی $+$ بهینه‌ساز کلاسیک
			\\
			\bottomrule
		\end{tabularx}
\end{table}

یک ابزار می‌تواند در چند الگوریتم ظاهر شود. \lr{Trotter} هم دینامیک می‌سازد، هم مسیر آدیاباتیک دیجیتال، هم ورودی \lr{QPE}.

\section{ارتباط}

مدل اجرا در فصل~\ref{chap:qmodel}، مدل‌های فیزیکی در فصل~\ref{chap:models}، رهیافت‌ها در فصل~\ref{chap:evo-adi}، و ابزارها در فصل‌های~\ref{chap:digital-tools} و~\ref{chap:spectral-tools} آمده‌اند.

\section{چه وقت استفاده می‌شود}

قبل از خواندن هر مقاله بپرسید: این کار مسئله است، الگوریتم است، یا ابزار؟ اگر هر سه را در یک جمله قاطی کردید، به این فصل برگردید.

\begin{remark}[جای توسعه]
تمایز شبیه‌سازی آنالوگ فاینمن در برابر شبیه‌سازی دیجیتال لوید؛ جایگاه نمونه‌برداری بوزونی به‌عنوان مسئله‌ای مجزا از شبیه‌سازی هامیلتونی.
\end{remark}
